\begin{textsample}{POS Dim 2 – unprofiled\_gpt – Score 7.00 – unprofiled\_gpt/t000939\_gpt.txt}  \label{ex:f2_pos_029}
What I ’m really getting at is that these moving image case studies are valuable precisely because they ’re \textbf{real} and co‑created, not something we ’ve just invented on paper.
They ’re grounded in authentic situations, with actual people collaborating to solve a \textbf{real} problem, and that authenticity gives them weight as learning tools.
We can package these as complete case-study units that include the video narrative plus integrated teaching and learning resources—things tutors can pick up and use in tutorials or \textbf{build} into online units.
The \textbf{idea} is that, say, if the theme is collaboration, the case offers a concrete, lived example of collaboration in action that both instructors and learners can work with effectively.
Higher education institutions are the obvious customers.
They struggle to keep up with new themes and emerging practices, and they often don't have the skills or the time to produce this kind of material \textbf{themselves}, especially not at the level of narrative quality we ’re talking about.
Their online content just doesn't move as fast as the kinds of stories we can generate.
Because this is real-life material, co-created with participants and designed from the outset as a learning resource, it ’s \textbf{far} \textbf{more} engaging—particularly for online students, who need something vivid to connect with.
That ’s why we can sell it at a premium : it ’s authentic, it ’s current, and it fills a gap universities can't easily fill on their \textbf{own}.

% matched lemmas: build, far, idea, more, own, real, themselves
\end{textsample}
